\documentclass[11pt,a4paper]{article}
\usepackage{amsmath,amssymb,amsthm,amsfonts}
\usepackage{geometry}
\geometry{margin=1.1in}
\usepackage{enumitem}
\usepackage{booktabs}
\usepackage{mathtools}
\usepackage{hyperref}
\usepackage[nameinlink,noabbrev]{cleveref}

\newtheorem{theorem}{Theorem}[section]
\newtheorem{proposition}[theorem]{Proposition}
\newtheorem{lemma}[theorem]{Lemma}
\newtheorem{corollary}[theorem]{Corollary}
\theoremstyle{definition}
\newtheorem{definition}[theorem]{Definition}
\newtheorem{example}[theorem]{Example}
\theoremstyle{remark}
\newtheorem{remark}[theorem]{Remark}

\newcommand{\Jnorm}{J_{\rm norm}}
\newcommand{\Jbound}{\frac{3\pi^2}{8}}

\title{Dual Spacetime 4-Valued Logic (D4L):\\
Amplitudes, Four States, Geometric Interference,\\
and Dual Hilbert Space}

\author{https://github.com/hypernumbernet}
\date{\today}

\begin{document}
\maketitle

\begin{abstract}
We fix Dual Spacetime 4-Valued Logic (D4L) as the logic of Dual Spacetime
Theory. A proposition is an admissible torsion configuration, equivalently
its PGA bivector \(\Omega\) and rotor \(R=\exp\Omega\). Measurement is the
normalised height \(\Jnorm\in[-1,1]\); collapse yields four labels
\(\lvert T\rangle,\lvert U\rangle,\lvert F\rangle,\lvert B\rangle\).
Scalar connectives are \(\min/\max/-\). Commutative overlap is the
parameter Killing form; non-commutative interference is the geometric
commutator of torsion bivectors. Two preorders on \([-1,1]\) separate
height from saturation. The dual sector carries a Hilbert space \(\mathbb{C}^2\): cyclic
generators obey the quaternion table, the flipped dual pairing is
Euclidean on \(\mathbb{R}^3\), and \(-i\sigma_a\) represents those
generators. For every dual rapidity the associated matrix rotor lies in
\(\mathrm{SU}(2)\). The closed-subspace lattice of \(\mathbb{C}^2\) is
orthomodular and non-distributive. The Clifford algebra \(\mathrm{Cl}(3,1)\)
acts on a Dirac spinor \(\mathbb{C}^4\) by matrices satisfying
\(\{\gamma^\mu,\gamma^\nu\}=2\eta^{\mu\nu}\); this representation is not
identified with the dual \(\mathbb{C}^2\). Amplitude geometry and the
Hilbert layer are two layers of one logic, related by an internal dictionary;
their operations are not identified. In particular the Killing form is
indefinite, self-overlap is not a Born probability, and the four labels
are not a PVM of \(\mathbb{C}^2\). Theorems below are machine-checked.
\end{abstract}

\section{Position}
\label{sec:position}

Dual Spacetime Theory encodes every massive configuration by six
rapidities \((\alpha_a,\beta_a)_{a=1,2,3}\): a usual-sector boost and a
dual-sector rotation on each spatial axis \cite{DST2025,Diophantine2026}.
The already-proved invariant of that pair is a signed height. D4L is the
logic of that geometry, in two layers.

The amplitude layer takes admissible configurations as propositions,
\(\Jnorm\) as the observable, and \(\min/\max/-\) as scalar connectives.
Killing overlap and the bivector commutator belong here: they act on
torsion parameters already present in the PGA core.

The dual Hilbert layer is the kinematics of the dual sector: a Euclidean
pairing on \(\mathbb{R}^3\), the quaternion table of the cyclic
generators, and their representation on \(\mathbb{C}^2\). The
closed-subspace lattice is the Birkhoff--von Neumann logic of a
two-dimensional spinor \cite{BvN1936}.

An internal dictionary relates the layers. Their operations are not
identified: the Killing form is not the spinor inner product,
\(\min/\max\) is not the meet and join of \(L(\mathbb{C}^2)\), and the
four labels are not a PVM. There is no \(\hbar\) and no derivation of
the Born rule from \(J\).

\section{Torsion and the admissible cone}
\label{sec:torsion}

\begin{definition}[Torsion scalar]
\label{def:J}
For a rotor pair write \(\Omega=R_{\rm usual}^\dagger R_{\rm dual}\) and
let \(\Omega_{\rm biv}\) be a principal logarithm. The parameter form used
throughout this paper is
\begin{equation}
\label{eq:J-raw}
J = \frac12\sum_{a=1}^3\bigl(\alpha_a^2-\beta_a^2\bigr).
\end{equation}
The combination \(\frac1{16}B(\Omega_{\rm biv},\Omega_{\rm biv})\) appearing
in older drafts differs from \eqref{eq:J-raw} by a constant factor under
the standard generator expansion; boundedness statements here always
refer to \(J\) and \(\Jnorm\) \cite{Diophantine2026}.
\end{definition}

\begin{definition}[Admissible configuration]
\label{def:admissible}
A configuration is \emph{admissible} when, on every axis,
\[
\alpha_a\ge 0,\qquad \beta_a\ge 0,\qquad \alpha_a+\beta_a\le \frac{\pi}{2}.
\]
\end{definition}

The cutoff \(\pi/2\) is the anti-synchronisation bound of the dual map:
usual and dual rapidities on one axis are not allowed to add past a right
angle.

\begin{theorem}[Torsion bound]
\label{thm:torsion-bound}
On the admissible cone of Definition~\ref{def:admissible},
\[
|J|\le\Jbound,\qquad |\Jnorm|\le 1,
\]
where
\begin{equation}
\label{eq:Jnorm}
\Jnorm := \frac{8}{3\pi^2}\,J.
\end{equation}
Equality \(|\Jnorm|=1\) holds if and only if every axis is pure boost, or
every axis is pure rotation. The image of \(\Jnorm\) on the admissible cone
is the whole interval \([-1,1]\). The bound is false if one keeps only the
principal-branch condition \(|\alpha_a+\beta_a|\le\pi/2\) and allows
negative rapidities.
\end{theorem}

For a compactification \(N\) the discrete bound is sharper
\cite{Diophantine2026}:
\[
|\Jnorm|\le\Bigl(\frac{4\lfloor N/4\rfloor}{N}\Bigr)^2.
\]
When \(4\mid N\) the right-hand side is \(1\) and both signs are attained.
When \(4\nmid N\) every admissible point satisfies the strict inequality
\(|\Jnorm|<1\).

\section{Amplitudes}
\label{sec:amp}

\begin{definition}[Amplitude]
\label{def:amp}
An \emph{amplitude} is an admissible continuous torsion configuration
\(p=(\alpha_a,\beta_a)\). Write
\[
\Omega(p)=\sum_{a=1}^3\Bigl(\frac{\alpha_a}{2}\,B^+_a+\frac{\beta_a}{2}\,B^-_a\Bigr),
\qquad
R(p)=\exp\Omega(p).
\]
Measurement and collapse are
\[
\langle p\rangle:=\Jnorm(p),\qquad
\mathrm{collapse}(p)
\in\{\lvert T\rangle,\lvert U\rangle,\lvert F\rangle,\lvert B\rangle\}.
\]
The adjoint is the usual--dual swap \(p^\dagger\), exchanging
\((\alpha_a)\) with \((\beta_a)\).
\end{definition}

\begin{theorem}[Adjoint]
\label{thm:adjoint}
The swap is an involution and flips the observable:
\( (p^\dagger)^\dagger=p \) and \(\langle p^\dagger\rangle=-\langle p\rangle\).
Every label is realised by some amplitude.
\end{theorem}

The four names are not a function of each other under adjoint: deepest
\(\lvert B\rangle\) maps to \(\lvert F\rangle\), interior \(\lvert B\rangle\)
maps to \(\lvert U\rangle\).

\section{Four states}
\label{sec:states}

\begin{definition}[D4L states]
\label{def:states}
On \(\Jnorm\in[-1,1]\),
\begin{align*}
\lvert T\rangle &\colon \Jnorm = 0
  &&\text{perfect usual--dual balance,}\\
\lvert U\rangle &\colon 0 < \Jnorm < 1
  &&\text{boost-dominant, not saturated,}\\
\lvert F\rangle &\colon \Jnorm = +1
  &&\text{every axis pure boost,}\\
\lvert B\rangle &\colon -1\le \Jnorm < 0
  &&\text{rotation-dominant, including the}\\
&&& \text{deepest point \(\Jnorm=-1\).}
\end{align*}
\end{definition}

The positive extreme is a singleton named False; the negative extreme is
kept inside Both. That is a logical convention, not a gap in the cone:
both extremes exist and are unique. Mixed extrema --- one axis boost,
another rotation --- never reach \(|\Jnorm|=1\). Saturation requires every
axis to vote the same way.

\begin{theorem}[Extremal labels]
\label{thm:extremal}
On an admissible configuration, \(\mathrm{collapse}(p)=\lvert F\rangle\)
if and only if \(p\) is purely hyperbolic, and \(\Jnorm(p)=-1\) if and
only if \(p\) is purely elliptic.
\end{theorem}

Interior \(\lvert B\rangle\) is a range, not a single fixed point. The
deepest point of that range is the wall \(\Jnorm=-1\).

Label \(\lvert T\rangle\) is likewise not a single geometric point.
The signed height \(\Jnorm=0\) holds both at the origin of the cone
and on a positive-mass balanced configuration \(\alpha_a=\beta_a\).
Those two are distinguished by a second observable, not by a fifth
name.

\begin{definition}[Mass]
\label{def:mass}
The unsigned torsional mass is
\[
M(p)=\frac12\sum_{a=1}^3\bigl(\alpha_a^2+\beta_a^2\bigr),\qquad
M_{\mathrm{norm}}=\frac{8}{3\pi^2}\,M.
\]
On the admissible cone \(0\le M_{\mathrm{norm}}\le 1\). Usual--dual
swap preserves \(M\). Uniform scaling multiplies \(M\) by \(k^2\).
\end{definition}

\begin{definition}[Vacuum and balanced massive]
\label{def:vacuum}
These are predicates on amplitudes, not extra D4L labels.
\begin{align*}
\mathrm{IsVacuum}(p)
  &\iff \mathrm{collapse}(p)=\lvert T\rangle \wedge M(p)=0,\\
\mathrm{IsBalancedMassive}(p)
  &\iff \mathrm{collapse}(p)=\lvert T\rangle \wedge 0<M(p).
\end{align*}
\(\mathrm{IsVacuum}\) is equivalent to all six rapidities vanishing,
and implies \(\Jnorm=0\). The converse is false.
\end{definition}

\begin{theorem}[Zero height is not vacuum]
\label{thm:T-split}
There exists an admissible configuration with \(\Jnorm=0\) and
\(M>0\). Consequently the paper step ``\(\Jnorm=0\) implies trivial
bases'' is false. Vacuum and balanced massive both inhabit
\(\lvert T\rangle\).
\end{theorem}

\section{Connectives}
\label{sec:conn}

\begin{definition}[Logical operations]
\label{def:ops}
On the signed height,
\[
\neg j = -j,\qquad
j\wedge k = \min(j,k),\qquad
j\vee k = \max(j,k).
\]
Negation is the usual--dual swap
\((\alpha_a,\beta_a)\mapsto(\beta_a,\alpha_a)\), which sends
\(J\mapsto -J\) and therefore \(\Jnorm\mapsto-\Jnorm\).
\end{definition}

\begin{theorem}[Non-explosion]
\label{thm:nexplosion}
For every \(j\in[-1,1]\),
\[
j\wedge\neg j = -|j|\le 0.
\]
Hence \(p\wedge\neg p\) is \(\lvert T\rangle\) or \(\lvert B\rangle\), and
never saturates \(\lvert F\rangle\). De Morgan duality and distributivity
hold as they do for \(\min/\max\) on \(\mathbb{R}\).
\end{theorem}

\section{Two orders, not Belnap FOUR}
\label{sec:order}

\begin{definition}
On \(\mathbb{R}\), and then by restriction on \([-1,1]\),
\[
j\preceq_\tau k \iff j\le k,
\qquad
j\preceq_\iota k \iff \lvert j\rvert\le\lvert k\rvert.
\]
The first is the height (boost-versus-rotation) order. The second is
the information (saturation) order.
\end{definition}

\begin{theorem}[Information bottom and tops]
\label{thm:info}
On \([-1,1]\), \(j=0\) is the unique information bottom, and the
information tops are exactly the walls \(j=\pm 1\). Equivalently, the
information bottom is the label \(\lvert T\rangle\), and the two tops
are \(\lvert F\rangle\) and deepest \(\lvert B\rangle\).
\end{theorem}

Interior \(\lvert B\rangle\) is strictly below deepest \(\lvert B\rangle\)
in the information order. The labels do not match Belnap's
\(\{\bot,t,f,\top\}\); no bilattice isomorphism is claimed.

\section{Geometric operations}
\label{sec:geo}

\begin{definition}[Overlap and interference]
\label{def:geo}
\[
\langle p\mid q\rangle := 8\sum_{a=1}^3\bigl(\alpha_a\alpha'_a-\beta_a\beta'_a\bigr),
\qquad
[\Omega(p),\Omega(q)] := \Omega(p)\Omega(q)-\Omega(q)\Omega(p).
\]
Rotor composition \(R(p)R(q)\) is a product in the PGA. It is not
pulled back through the Baker--Campbell--Hausdorff formula to a third
torsion configuration.
\end{definition}

\begin{theorem}[Pairing and non-commutativity]
\label{thm:geo}
The overlap is symmetric, and \(\langle p\mid p\rangle=16\,J(p)\),
equivalently a constant multiple of \(\Jnorm(p)\). Distinct-axis
torsion bivectors need not commute: if \(p\) is a pure boost on axis
\(0\) and \(q\) a pure rotation on axis \(1\), then
\([\Omega(p),\Omega(q)]\ne 0\). Same-axis hyperbolic and cyclic
generators do commute, so the corresponding one-axis interference
vanishes.
\end{theorem}

\begin{theorem}[Indefinite pairing]
\label{thm:indef}
There exist configurations with \(\langle p\mid p\rangle>0\) and with
\(\langle q\mid q\rangle<0\). In particular the pairing is not
positive definite.
\end{theorem}

The witnesses are the pure-boost and pure-rotation extremals already
used for \(\lvert\Jnorm\rvert=1\).

\begin{theorem}[Not a Born probability]
\label{thm:born}
There exists \(p\) with \(\langle p\mid p\rangle<0\). A signed quantity
cannot be a Born probability.
\end{theorem}

The pairing is not the spinor inner product
(Section~\ref{sec:spinor}). Sandwich evaluation of rotors is
independent of the gravity chart; preservation of the degenerate
quadratic form is not claimed. Geometric composition is associative as
a PGA product and is not a lattice operation.

\section{Dual Hilbert layer}
\label{sec:hilbert}

The dual sector already contains a Euclidean three-space, a quaternion
table, and a spinor representation on \(\mathbb{C}^2\). This section
records that kinematics and the dictionary that relates it to the
amplitude layer. Spatial axes are indexed by \(a=0,1,2\).

\subsection{Dual and usual sectors}
\label{sec:sector}

\begin{definition}
\(p\) is dual-sector when \(\alpha=0\), usual-sector when \(\beta=0\).
Write \(p_\beta\) and \(p_\alpha\) for the corresponding embeddings of
Euclidean \(\mathbb{R}^3\).
\end{definition}

\begin{theorem}[Definite restrictions]
\label{thm:sector}
On the dual sector,
\[
\langle p_\beta\mid p_\gamma\rangle
=-8\langle\beta,\gamma\rangle_{\mathbb{R}^3},
\]
hence negative definite. On the usual sector the pairing is
\(+8\langle\alpha,\gamma\rangle_{\mathbb{R}^3}\), hence positive
definite.
\end{theorem}

The inner-product space is the linear span of each sector, not the
admissible cube \(\alpha_a,\beta_a\ge 0\), \(\alpha_a+\beta_a\le\pi/2\).
Superposition is a statement about that span.

\subsection{Quaternion table}
\label{sec:quat}

Write \(\Gamma_a\) for the cyclic generators of the dual sector. They
satisfy \(\Gamma_a^2=-1\) and the off-diagonal products below.

\begin{theorem}[Quaternion multiplication]
\label{thm:quat}
\[
\begin{aligned}
\Gamma_0\Gamma_1&=\Gamma_2,\\
\Gamma_1\Gamma_2&=\Gamma_0,\\
\Gamma_2\Gamma_0&=\Gamma_1,
\end{aligned}
\qquad
\begin{aligned}
\Gamma_1\Gamma_0&=-\Gamma_2,\\
\Gamma_2\Gamma_1&=-\Gamma_0,\\
\Gamma_0\Gamma_2&=-\Gamma_1.
\end{aligned}
\]
\end{theorem}

The identities follow from the vector anticommutators in \(G(3,1,1)\).
For a pure dual configuration
\[
\Omega(p_\beta)=\sum_a(\beta_a/2)\,\Gamma_a,
\]
and on axis \(0\) this collapses to \((\theta/2)\,\Gamma_0\).

\subsection{Spinor Hilbert space}
\label{sec:spinor}

\begin{definition}[Dual spinor]
The dual spinor space is \(\mathbb{C}^2\) with the standard Hermitian
product. Pauli matrices \(\sigma_x,\sigma_y,\sigma_z\) are the usual
ones. The representation is \(\rho(\Gamma_a):=-i\sigma_a\).
That sign makes \(IJ=K\) agree with the cyclic products in \(G(3,1,1)\)
and sends the dual rotor to the standard \(\mathrm{SU}(2)\) factor
\(\exp(-i\beta\cdot\sigma/2)\).
\end{definition}

\begin{theorem}[Representation]
\label{thm:rep}
\(\rho(\Gamma_a)^2=-1\) and
\(\rho(\Gamma_0)\rho(\Gamma_1)=\rho(\Gamma_2)\).
Each \(\rho(\Gamma_a)\) is anti-Hermitian.
\end{theorem}

The dual rotor matrix is
\[
U(\beta)
=\exp\Bigl(\sum_a(\beta_a/2)(-i\sigma_a)\Bigr)
=\exp(-i\beta\cdot\sigma/2).
\]
On a single axis this is the Rodrigues formula
\(\cos(\theta/2)\,I+\sin(\theta/2)\,(-i\sigma_a)\). For a general dual
rapidity the same trigonometric form holds along the unit cyclic
combination determined by \(\beta\).

\begin{theorem}[Dual rotor in \(\mathrm{SU}(2)\)]
\label{thm:su2}
For every dual rapidity \(\beta\in\mathbb{R}^3\), the matrix \(U(\beta)\)
is unitary of determinant \(1\). Thus \(U(\beta)\in\mathrm{SU}(2)\).
\end{theorem}

A general identification of the Banach exponential in \(G(3,1,1)\) with
this matrix exponential is not claimed; that would require control of
the Baker--Campbell--Hausdorff formula for mixed generators.

\(\dim_{\mathbb{C}}\mathbb{C}^2=2\), so a projective measurement
has at most two orthogonal outcomes.

\subsection[Quantum logic of C2]{Quantum logic of \(\mathbb{C}^2\)}
\label{sec:ql}

\begin{definition}
A quantum proposition is a linear subspace of \(\mathbb{C}^2\).
Meet is intersection, join is the sum, negation is the orthogonal
complement. Finite dimension makes every subspace closed.
\end{definition}

\begin{theorem}[Orthomodular law]
\label{thm:om}
If \(A\le B\) then \(A\vee(A^\perp\wedge B)=B\).
\end{theorem}

This is the orthogonal decomposition of a finite-dimensional subspace.

\begin{theorem}[Non-distributivity]
\label{thm:ndist}
Let \(e_0=(1,0)\), \(e_1=(0,1)\), \(d=(1,1)\), and write \(\ell_v\) for
\(\mathbb{C}v\). Then
\[
\ell_{e_0}\wedge(\ell_{e_1}\vee\ell_d)
\neq
(\ell_{e_0}\wedge\ell_{e_1})\vee(\ell_{e_0}\wedge\ell_d).
\]
The left side is \(\ell_{e_0}\); the right side is \(\{0\}\).
\end{theorem}

The scalar connectives of Section~\ref{sec:conn} distribute. A
distributive lattice cannot be isomorphic to the subspace lattice of a
Hilbert space of dimension at least two.

\begin{theorem}[At most two orthogonal atoms]
\label{thm:two}
There do not exist three pairwise orthogonal one-dimensional subspaces
of \(\mathbb{C}^2\).
\end{theorem}

\subsection{Projectors}
\label{sec:proj}

A primitive idempotent in the algebra must be built from generators
that square to \(+1\). If \(g^2=+1\), then \((1+g)/2\) is idempotent.
Write \(j:=e_0e_1\) for the hyperbolic generator of axis \(0\) (so
\(j^2=+1\)). This \(j\) is not the time vector \(\gamma^0=e_0\), which
squares to \(-1\) in signature \((-,+,+,+)\). The working projectors
are
\[
P_{\mathrm{spin}}=\frac{1+e_0e_1}{2},\qquad
P_L=\frac{1-e_1}{2},\qquad
P_R=\frac{1+e_1}{2}.
\]
They satisfy \(P_L^2=P_L\), \(P_R^2=P_R\), \(P_{\mathrm{spin}}^2=P_{\mathrm{spin}}\),
and \(P_L+P_R=1\).

\begin{theorem}[Draft projectors fail]
\label{thm:proj-fail}
The formulae \((1\pm i)/2\) built from the \(\mathrm{Cl}(3,1)\)
pseudoscalar \(i\) with \(i^2=-1\) are not idempotent:
\(\bigl((1-i)/2\bigr)^2=-i/2\). The product \(P_{\mathrm{spin}}P_R\) is
likewise not idempotent: \(e_0e_1\) and \(e_1\) anticommute.
\end{theorem}

\begin{theorem}[Commuting projector]
\label{thm:proj-ok}
The generators \(e_1\) and \(e_0e_2\) commute and both square to
\(+1\). Their product projector
\[
P=\frac{1+e_1}{2}\cdot\frac{1+e_0e_2}{2}
\]
is idempotent and nonzero. The principal left ideal it generates is
nonzero and has real dimension at least two. Irreducibility of that
ideal as a spinor module, and a complex dimension of four, are not
claimed.
\end{theorem}

\subsection{Clifford relations and the Dirac spinor}
\label{sec:dirac}

The Dirac matrices of signature \((-,+,+,+)\) are the grade-1 generators
\(\gamma^\mu\) of \(\mathrm{Cl}(3,1)\). They obey
\[
\{\gamma^\mu,\gamma^\nu\}=2\eta^{\mu\nu},\qquad
\eta=\mathrm{diag}(-1,+1,+1,+1).
\]
In particular \((\gamma^0)^2=-1\), so \(\gamma^0\neq j\).

\begin{theorem}[Dirac representation]
\label{thm:dirac-rep}
There exist matrices \(\gamma^\mu\in M_4(\mathbb{C})\) satisfying the
same Clifford relations. They act on a Dirac spinor space
\(\mathbb{C}^4\). This space has complex dimension four, hence is not
the dual spinor \(\mathbb{C}^2\). As vector spaces one has a Weyl
splitting \(\mathbb{C}^4\simeq\mathbb{C}^2\oplus\mathbb{C}^2\); that
splitting is a dictionary, not an identification of the dual-sector
kinematics with a chiral Dirac component.
\end{theorem}

The Dirac equation, a mass term, and covariance under the full rotor
group are not derived here.

\subsection{Internal dictionary}
\label{sec:dict}

\begin{center}
\begin{tabular}{@{}ll@{}}
\toprule
Amplitude layer & Dual Hilbert layer \\
\midrule
pure dual amplitude & \(\beta\in\mathbb{R}^3\) and \(U(\beta)\in\mathrm{SU}(2)\) \\
usual--dual swap & leaves the dual sector; not a Hilbert adjoint \\
Killing overlap & distinct from the spinor inner product \\
vanishing interference on one axis & commuting generators (necessary only) \\
\(\Jnorm\) four-state collapse & coarse-graining of a continuous scalar \\
\(\min/\max\) & not the meet/join of \(L(\mathbb{C}^2)\) \\
dual \(\mathbb{C}^2\) & not the Dirac \(\mathbb{C}^4\) \\
\bottomrule
\end{tabular}
\end{center}

\begin{theorem}[Four labels are not a PVM]
\label{thm:pvm}
There is no assignment of one-dimensional subspaces of
\(\mathbb{C}^2\) to the four D4L labels that makes distinct
labels pairwise orthogonal.
\end{theorem}

\begin{theorem}[Compatibility is not same-axis]
\label{thm:compat}
If two configurations are supported on a single common spatial axis,
their interference vanishes. The converse is false: vanishing
interference does not force a common single axis. A two-axis boost,
for instance, has vanishing self-commutator and is not single-axis.
\end{theorem}

\section{Designated values and consequence}
\label{sec:cons}

The four labels are not yet a logic until one says which heights
\emph{hold}. D4L uses two designated predicates on \(\Jnorm\in[-1,1]\),
not one.

\begin{definition}[Designated values]
\label{def:desig}
\[
\mathrm{HoldsT}(j)\iff j=0,
\qquad
\mathrm{HoldsNotF}(j)\iff j<1.
\]
The first is synchrony (the label \(\lvert T\rangle\)). The second is
non-refutation (every label except \(\lvert F\rangle\)).
\end{definition}

Formulas are generated by atoms, \(\neg\), \(\wedge\), and \(\vee\),
and are evaluated by the scalar operations \(\neg\), \(\wedge\), \(\vee\)
already fixed on heights.
There is no implication connective on the amplitude heights: the
residuum of \(\min\) would take its top at \(+1=\lvert F\rangle\).
Implication lives on the discrete proof-status layer
(Section~\ref{sec:regime}).

\begin{definition}[Consequence]
\label{def:entail}
A valuation assigns each atom a height in \([-1,1]\). \(\Gamma\models_T\psi\)
when every valuation that synchronises \(\Gamma\) synchronises \(\psi\).
\(\Gamma\models_{\lnot F}\psi\) is the same clause with non-refutation.
\end{definition}

Two-valued \emph{restrictions} of the interval are not unique, because
D4L's names \(\lvert T\rangle,\lvert F\rangle\) are not classical
true/false.

\begin{definition}[Two fragments]
\label{def:frag}
The \emph{named} fragment is \(\{0,1\}=\{\lvert T\rangle,\lvert F\rangle\}\).
The \emph{wall} fragment is \(\{-1,+1\}=\{\text{deepest }\lvert B\rangle,\lvert F\rangle\}\).
\end{definition}

\begin{theorem}[Named fragment is not Boolean]
\label{thm:named-not-bool}
\(\neg 0=0\), so the named fragment contains a negation fixed point.
It is not closed under \(\neg\): \(\neg 1=-1\) leaves \(\{0,1\}\).
\end{theorem}

\begin{theorem}[Walls swap]
\label{thm:wall-swap}
If \(j\in\{-1,+1\}\) then \(\neg j\in\{-1,+1\}\) and \(j\neq\neg j\).
\end{theorem}

The wall fragment is the faithful classical core. The named fragment
matches the labels and is already not classical.

\section{Examples two-valued logic cannot host}
\label{sec:examples}

These are changes of semantics, not new theorems of Peano arithmetic.
Unconditional FLT / Beal and a refutation of G\"odel's theorems are
not claimed.

\subsection{Negation fixed point}
\label{sec:ex-fp}

\begin{theorem}[Fixed point of negation]
\label{thm:ex-fp}
The condition \(j=\neg j\) is equivalent to \(j=0\). It is
unsatisfiable on the wall fragment and uniquely realised in D4L by
\(\lvert T\rangle\), including by an amplitude.
\end{theorem}

Classically \(P\leftrightarrow\neg P\) has no model. D4L hosts it at
perfect usual--dual balance. This is not the G\"odel sentence.

\subsection{Contradiction without explosion}
\label{sec:ex-efq}

\begin{theorem}[Boolean EFQ is vacuous]
\label{thm:bool-efq}
On \(\mathrm{Bool}\), \(P\) and \(\neg P\) cannot both be
\(\mathtt{true}\). On the walls, \(j\) and \(\neg j\) cannot both be
\(+1\).
\end{theorem}

\begin{theorem}[D4L does not explode]
\label{thm:d4l-efq}
Let \(P\) be atom \(0\) and \(Q\) atom \(1\). Then
\(\{P,\neg P\}\not\models_T Q\) and
\(\{P,\neg P\}\not\models_{\lnot F} Q\). A witness sends \(P\) to \(0\)
and \(Q\) to \(1\). A second witness sends \(P\) to \(-\tfrac12\)
(interior \(\lvert B\rangle\)) and \(Q\) to \(1\). The conjunction
\(j\wedge\neg j\) is always non-refuted.
\end{theorem}

An inconsistent pair of observations therefore does not force an
unrelated atom onto the positive wall.

\subsection{\texorpdfstring{The constraint \(\Jnorm<1\)}{The constraint Jnorm < 1}}
\label{sec:ex-notF}

\begin{theorem}[Named remainder is a singleton]
\label{thm:named-notF}
On the named fragment, \(\mathrm{HoldsNotF}\) collapses to
\(\mathrm{HoldsT}\). On D4L the same constraint is inhabited by
\(\lvert T\rangle\), \(\lvert U\rangle\), and \(\lvert B\rangle\).
\end{theorem}

\subsection{Complementarity}
\label{sec:ex-comp}

\begin{theorem}[No Boolean assignment of complementary rays]
\label{thm:ex-comp}
The subspace lattice of \(\mathbb{C}^2\) is not distributive, hence not
a Boolean algebra. It has at most two pairwise orthogonal atoms. The
four D4L labels are not a PVM. The distributive scalar connectives
cannot present that lattice.
\end{theorem}

\section{Proof-status layer}
\label{sec:regime}

A DST proof artefact --- a core lemma, a live bridge, a refuted
diagnostic, a bookkeeping identity --- is not a torsion configuration.
It carries a discrete status that reuses the four names, but the
observable is not \(\Jnorm\). On this carrier negation \emph{is} a
function of the four labels, and implication is a table, not the
residuum of \(\min\).

\begin{definition}[Regime operations]
\label{def:regime}
Negation and the lattice operations follow the amplitude
representatives \(\lvert T\rangle\mapsto 0\), \(\lvert U\rangle\mapsto 1/2\),
\(\lvert F\rangle\mapsto 1\), \(\lvert B\rangle\mapsto -1/2\). In particular
\(F\wedge\neg F=B\). A second operation, the establishedness meet,
packages two artefacts along \(F<B<U<T\), so that a
proved core with a live bridge is still live: the meet of \(T\) and
\(U\) is \(U\). Implication is
\[
\begin{aligned}
T\to\phi&=\phi,\\
F\to\phi&=T,\\
U\to T&=U,\\
B\to T&=U.
\end{aligned}
\]
The last two clauses are the point: an open or conflicted programme
does not establish a theorem. No Belnap-FOUR isomorphism is claimed.
\end{definition}

Designated predicates on statuses are the same two names:
\(\mathrm{HoldsT}_R\) means status \(T\);
\(\mathrm{HoldsNotF}_R\) means status \(\ne F\).

Additive motors and multiplicative amplitudes occupy distinct seats at
the same label \(\lvert T\rangle\). Pure translation has vanishing
torsion and is vacuum. A balanced massive configuration is not vacuum.
The mixed product \(R\cdot T\) is not identified with
\(\exp(\Omega_{\mathrm{biv}})\) when the factors fail to commute. The
balanced residual class excludes window seeds: a half-window
configuration is not a balanced residual.

\section{Discrete amplitudes and amplification}
\label{sec:discrete-dyn}

On the finite torus, when \(4\nmid N\) the sharp discrete bound is
strict, so no admissible point saturates \(\lvert F\rangle\). When
\(4\mid N\) both walls are attained.

Pure-boost scaling multiplies \(\Jnorm\) by \(k^2\). Synchrony is a
fixed point. A label \(\lvert U\rangle\) can be driven onto
\(\lvert F\rangle\), or out of the admissible cone, where collapse is
undefined. That is the logical reading of the shared no-go; it is not
new arithmetic.

Vacuum is a fixed point of every admissible scaling. A small
balanced-massive seed stays at \(\lvert T\rangle\) with mass
\(k^2 M\) and may remain admissible; a large balanced seed leaves
the cone. That cone exit is invisible to the signed-height
classifier: \(\Jnorm\) stays \(0\), so the four-state label is still
\(\lvert T\rangle\). Mass and the cone predicate are the
observables that see it. This is why a Beal balanced seed
\(\theta=\log 2/m\) can survive \(m\)-fold scaling inside the cone.

Height and modular winding are two measurements. Real-scale
admissibility and nonzero winding cannot hold together. That
incompatibility is a separation of observables, not a Both-status on
one observable.

\subsection{Diophantine regimes}
\label{sec:ex-regime}

\begin{theorem}[Three-layer split]
\label{thm:ex-regime}
Let atom \(0\) be a proved core, atom \(1\) a diagnostic, atom \(2\) a
live bridge, and atom \(3\) a classical conjecture. Named two-valued
logic \(\{T,F\}\) cannot host the live status \(U\). The wall fragment
\(\{B,F\}\) cannot host the core \(T\). D4L realises the assignment
core \(T\), diagnostic \(F\), live \(U\), conjecture \(U\). The core
alone does not \(T\)-entail the conjecture. A window \(T\) together
with an ill-posed NoGo \(B\) does not force the conjecture off \(F\).
The implication \(U\to T\) is not designated.
\end{theorem}

The same algebra hosts a Beal residual atlas: closed slices (including
the Fermat and Darmon--Merel positions, and the signature \((n,n,5)\))
sit at \(T\); refuted diagnostics sit at \(F\); bookkeeping identities
sit at \(B\); live residuals and the classical conjecture sit at \(U\).
The closed slices do not \(T\)-entail classical Beal. A dual-axis
Fermat atlas likewise places the additive core and the real \(L^p\)
dichotomy at \(T\), demoted single-axis modular and balanced-seed
diagnostics at \(F\), and the mixed-motor residual together with
classical FLT at \(U\). Closed Fermat slices do not \(T\)-entail
classical FLT.

This is why DST must not treat a single bound, or a refuted diagnostic,
as entailing the seven classical conjectures. The regime layer records
that prohibition. It does not prove any unproved bridge, and it does
not claim unconditional FLT, Beal, or abc.

\section{Self-reference}
\label{sec:godel}

Let \(G\) be the G\"odel sentence: ``\(G\) is not provable in the
system.'' In D4L, provability is read as a drive toward synchrony
\(\Jnorm\to 0\), i.e.\ toward \(\lvert T\rangle\). The honest constraint
on a sentence forbidden from sitting at the unique all-boost world is
exactly \(\mathrm{HoldsNotF}\) of Section~\ref{sec:ex-notF}:
\[
\Jnorm(G)<1.
\]
That constraint leaves \(G\) in \(\lvert T\rangle\), \(\lvert U\rangle\),
or \(\lvert B\rangle\) as labels, or at \(\lvert T\rangle\) if one
further demands synchrony. Conjunction does not explode
(Theorem~\ref{thm:d4l-efq}).

This is a change of semantics, not a proof that G\"odel's theorems fail
in classical arithmetic. Completeness talk over the negative-torsion
sector is likewise a reading of admissible rotors, not a completion of
the Hilbert programme inside first-order Peano arithmetic. The theorems
of Section~\ref{sec:examples} are the fixed content; the
reading in this section is not a new incompleteness theorem.

\section{What D4L does and does not show}
\label{sec:claims}

\begin{itemize}[leftmargin=1.4em]
  \item It \emph{does} take admissible torsion configurations as the
        carriers of propositions, with \(\Jnorm\) as the observable.
  \item It \emph{does} partition \([-1,1]\) into four named states, with
        both walls attained on the admissible cone.
  \item It \emph{does} make negation geometric: the usual--dual swap.
  \item It \emph{does} exhibit a commutative pairing and a
        non-commutative interference term already proved in the
        algebra.
  \item It \emph{does} separate height from saturation by two
        preorders on \([-1,1]\).
  \item It \emph{does} keep \(p\wedge\neg p\) out of \(\lvert F\rangle\).
  \item It \emph{does} distinguish two designated predicates
        (\(\mathrm{HoldsT}\), \(\mathrm{HoldsNotF}\)) and two
        two-valued fragments (named versus walls).
  \item It \emph{does} host \(P=\neg P\) uniquely at \(\lvert T\rangle\),
        refuse explosion of \(\{P,\neg P\}\), inhabit
        \(\Jnorm<1\) by three labels, and host the DST three-layer
        split on a discrete proof-status algebra.
  \item It \emph{does} put implication on that status algebra, with
        \(U\to T=U\), and refuse the residuum of \(\min\) on heights.
  \item It \emph{does} split label \(\lvert T\rangle\) by a second
        observable \(M\): vacuum versus balanced massive, without a
        fifth name.
  \item It \emph{does} forbid discrete \(\lvert F\rangle\) when
        \(4\nmid N\), and read \(k\)-fold amplification as a partial
        map on labels, including cone exit of balanced massive that
        the height classifier cannot see.
  \item It \emph{does} prove that the Killing pairing is indefinite and
        that self-overlap is not a Born probability.
  \item It \emph{does} extract a Euclidean dual sector and a quaternion
        table for the cyclic generators.
  \item It \emph{does} represent that table on \(\mathbb{C}^2\), place
        the dual rotor matrix in \(\mathrm{SU}(2)\) for every dual
        rapidity, and record the orthomodular, non-distributive lattice
        of closed subspaces.
  \item It \emph{does} reject the draft projectors \((1\pm i)/2\) and
        \(P_{\mathrm{spin}}P_R\), and exhibit a commuting square-\(+1\)
        pair whose product is idempotent, with a left ideal of real
        dimension at least two.
  \item It \emph{does} realise the Clifford relations of
        \(\mathrm{Cl}(3,1)\) on a Dirac spinor \(\mathbb{C}^4\), distinct
        from the dual \(\mathbb{C}^2\).
  \item It \emph{does} host Beal and Fermat residual atlases on the
        regime algebra, without \(T\)-entailing the classical
        conjectures.
  \item It does \emph{not} identify the scalar connectives with the
        meet and join of \(L(\mathbb{C}^2)\), nor the four labels with
        a PVM, nor vanishing interference with a common single axis.
  \item It does \emph{not} derive the Born rule from \(J\) or from
        Killing overlap.
  \item It does \emph{not} derive the Dirac equation, position or
        momentum operators, or a solution of the measurement problem.
  \item It does \emph{not} claim irreducibility of the commuting left
        ideal, nor a complex dimension of four for that ideal.
  \item It does \emph{not} refute G\"odel's theorems in classical
        arithmetic.
  \item It does \emph{not} claim unconditional Beal (or FLT, abc)
        from the \(T\)-split or from mass scaling. Balanced massive
        seeds may remain admissible after amplification.
  \item It does \emph{not} pull \(R(p)R(q)\) back to a torsion
        configuration when \([\Omega(p),\Omega(q)]\ne 0\).
  \item It does \emph{not} identify raw \(J\) with a value in
        \([-1,1]\). That interval belongs to \(\Jnorm\); the raw
        ceiling is \(\Jbound\).
\end{itemize}

The statements above are machine-checked.

\begin{thebibliography}{9}
\raggedright

\bibitem{DST2025}
Dual Spacetime Theory: Gravity as Torsion between Dual Spacetimes (2025),
arXiv:2512.xxxxx.

\bibitem{Diophantine2026}
Discrete Biquaternionic Double Spacetime: Rigorous Proofs of Diophantine
Conjectures via Torsional Mismatch and PGA Embedding,
\texttt{dst-diophantine.tex} (2026). Lean companion:
\url{https://github.com/hypernumbernet/DstDiophantine}.

\bibitem{BvN1936}
G.~Birkhoff and J.~von Neumann, The logic of quantum mechanics,
Ann.\ Math.\ 37 (1936).

\bibitem{Godel1931}
K.~G\"odel, \"Uber formal unentscheidbare S\"atze der Principia Mathematica
und verwandter Systeme~I, Monatshefte Math.\ Phys.\ 38 (1931).

\bibitem{Priest2006}
G.~Priest, \emph{In Contradiction: A Study of the Transconsistent},
Oxford University Press (2006).

\end{thebibliography}

\end{document}
